% =====================================================================================
%  ab-ovo-overview.pl.tex  --  POLSKA EDYCJA ab-ovo-overview.tex
%
%  TO NIE JEST ARTYKUL NAUKOWY I NIE POWSTAL POD STANDARDEM RESEARCH-DOCUMENTATION.
%
%  Powiedzenie tego tutaj, wprost i na gorze, jest sednem tego komentarza. Dokument jest
%  zlozony jak artykul, bo zlozony przeglad latwiej komus wreczyc niz katalog markdowna, a
%  czytelnik, ktory pomyli forme z gatunkiem, bedzie oczekiwal rzeczy, ktorych ten dokument
%  nie ma i nie moze deklarowac: pytania badawczego, metody, danych, wyniku, sekcji o
%  zagrozeniach dla trafnosci, przegladu prac pokrewnych ani odkrycia, ktore da sie cytowac.
%  Nie ma zadnej z nich. To jest PRZEGLAD PROJEKTU.
%
%  NIE WNOSI TEZ ZADNEGO FAKTU, ktorego nie niesie jego zrodlo w markdownie. Kazde zdanie
%  ponizej stoi w README.md, docs/architecture/00-ARCHITECTURE.md, docs/adr/ albo
%  docs/ux/UI-UX.md, i to te pliki sa zrodlem prawdy. Jesli ten papier i ktorys z nich sie
%  nie zgadzaja, racje ma markdown, a ten plik jest nieaktualny.
%
%  JEST SIOSTRZANA EDYCJA ANGIELSKA, ab-ovo-overview.tex, i to ona jest oryginalem. Ta
%  edycja nie jest niezaleznym tlumaczeniem: slownictwo jest dobrane tak, by zgadzalo sie z
%  tym, ktorym produkt juz mowi do polskiego czytelnika (web/app/src/lib/i18n/chrome.ts) i
%  z docs/START-HERE.pl.md. Inaczej jedno pojecie zyskaloby dwie polskie nazwy w jednym
%  majatku. Identyfikatory kodu - sciezki tras, nazwy klas, nazwy plikow - zostaja po
%  angielsku w obu edycjach, bo to sa prawdziwe nazwy.
%
%  DRYF, KTORY PRZYJMUJEMY. Kuratorska prezentacja dokumentu markdown nie jest
%  przeksztalceniem mechanicznym, wiec nic nie wymusza, by edycja jednego pliku dosiegla
%  drugiego. Zapisujemy to jako przyjeta konsekwencje, a nie udajemy, ze pokrywa to jakas
%  konwencja. Gdyby obie edycje rozjechaly sie widocznie wiecej niz raz, tanie mechaniczne
%  sprawdzenie (diff naglowkow sekcji) jest wlasciwa odpowiedzia, a nie pelna regula
%  sprzezenia zmian.
%
%  BUILD. .github/workflows/build-overview-pdf.yml renderuje ten plik i jego angielskie
%  rodzenstwo na workflow_dispatch i na niczym wiecej, i wrzuca oba PDF-y jako jeden
%  artefakt przebiegu. PDF NIGDY NIE JEST COMMITOWANY: zacommitowany PDF jest binarka w
%  kazdym diffie i druga kopia dokumentu, ktorego zrodlo juz tu jest.
%
%  SILNIK. W odroznieniu od edycji angielskiej ten plik ZAWIERA ZNAKI SPOZA ASCII, bo
%  polszczyzna bez diakrytykow to zla polszczyzna. Kompiluje sie pod pdflatex-em ze
%  standardowym TeX Live dzieki inputenc[utf8], fontenc[T1] i babel[polish]; babel jest tym,
%  co daje tej edycji poprawne dzielenie wyrazow i polskie konwencje typograficzne, a bez
%  niego tekst bylby przetlumaczony, ale wciaz skladany po angielsku.
%
%  JEDEN ZAKAZANY ZNAK. Babel[polish] czyni " znakiem aktywnym, wiec ten plik nie uzywa go
%  nigdzie w tresci; cudzyslow sklada sie jako ,, i ''. Ten komentarz jest jedynym
%  miejscem, gdzie o tym mowa, i najtansza wersja tej reguly.
%
%  RYSUNKI SA RENDEROWANE, A NIE RYSOWANE TUTAJ. Kazdy z nich to diagram Mermaida, ktorego
%  zrodlem jest docs/diagrams/<id>-<slug>.pl.mmd - to samo zrodlo, ktore docs/DIAGRAMS.pl.md
%  renderuje na GitHubie, trzymane identycznie przez scripts/check-diagrams.mjs. Uruchom
%  `node scripts/render-diagrams.mjs` PRZED latexmk-iem; workflow tak robi, a bez tego build
%  umiera na pierwszym \includegraphics.
% =====================================================================================

\documentclass[11pt,a4paper]{article}

\usepackage[T1]{fontenc}
\usepackage[utf8]{inputenc}
\usepackage[polish]{babel}
\usepackage[a4paper,margin=28mm]{geometry}
\usepackage{booktabs}
\usepackage{graphicx}
\usepackage{enumitem}
\usepackage{parskip}
\usepackage[hidelinks]{hyperref}

\newcommand{\code}[1]{\texttt{#1}}

\title{ab-ovo: platforma edukacyjna, która zamyka w sobie\\
\emph{Matematykę od zera dla inżyniera AI}}
\author{konradcinkusz}
\date{20 września 2026}

\begin{document}
\maketitle

\begin{abstract}
\noindent
ab-ovo to platforma edukacyjna zbudowana wokół książki złożonej z 47 programów nauczania
programowanego, po angielsku i po polsku, wraz z ćwiczeniami komputerowymi tej książki. Ten
dokument jest przeglądem projektu: czym system jest, czym celowo nie jest, co istnieje dziś w
repozytorium i co jest planowane. Nie jest artykułem naukowym i nie podaje żadnego wyniku.
Każde zdanie w nim pochodzi z dokumentacji markdown repozytorium, która jest źródłem prawdy;
gdy oba się nie zgadzają, rację ma markdown. Rysunki są własnymi diagramami repozytorium,
wyrenderowanymi z tych samych źródeł, które renderuje GitHub.
\end{abstract}

\section{Czym jest ab-ovo}

Książka, którą zamyka w sobie, stoi na mechanizmie, którego papier nie potrafi wymusić. Ramka
Strouda prosi czytelnika o coś \emph{zanim} cokolwiek mu powie, a kolejna ramka otwiera się
odpowiedzią, którą miał już zapisać. Czytelnik, który przelatuje wzrokiem, nie dostaje z tej
struktury nic, a drukowana strona nie ma jak tego zauważyć.

ab-ovo jest oprogramowaniem, które ma jak. Czytelnik czyta ramkę, zobowiązuje się do
odpowiedzi i odsłania kolejną ramkę, która otwiera się odpowiedzią. Odsłonięcie jest nawigacją,
a nie przełącznikiem, więc odpowiedź na ramkę widoczną na ekranie renderuje żądanie o kolejną
i nic wcześniej: nie ma jej w dokumencie, zamiast być w nim ukrytą, a czytelnik, który otworzy
inspektor, nie znajdzie niczego. Tam, gdzie ramka o coś prosi, pod nią stoi arkusz --- linia
odpowiedzi, notatnik liczący arytmetykę, płótno do szkicu --- a wszystkie trzy są opcjonalne i
lokalne dla przeglądarki czytelnika. Ćwiczenia komputerowe książki robi się w osobnym panelu
laboratorium, w którym Python działa we własnej przeglądarce czytelnika; opuściły one pętlę
lektury celowo, bo czytelnik książki matematycznej nie powinien musieć pisać kodu, żeby
odpowiedzieć na ramkę.

\begin{figure}[ht]
\centering
\includegraphics[width=0.78\linewidth]{../diagrams/rendered/b1-reader-loop.pl.pdf}
\caption{Pętla czytelnika. Każde pole działa bez konta i bez backendu; dwie kropkowane
krawędzie z niej wychodzące to jedyne miejsca, gdzie pojawia się serwer, i żadna nie leży na
ścieżce.}
\end{figure}

Pierwszy wymóg produktu ogranicza wszystko inne w tym dokumencie: \textbf{pętla czytelnika
działa bez konta i bez backendu.} Ramki są serwowane razem z witryną jako wersjonowana paczka
treści, ćwiczenia działają po stronie klienta, a nic, co czytelnik pisze na ramce, nie jest
przechowywane nigdzie poza jego własną przeglądarką. Konto dokłada synchronizację pozycji w
lekturze między maszynami i nic więcej.

\section{Czym celowo nie jest}

System zbiera wyniki, a wyniki o uczeniu się dryfują ku byciu wynikami o uczących się, bo to
jest łatwiejsza liczba do wyprodukowania i ta, która wygląda jak postęp. Odpowiedzią projektu
jest antycel wypowiedziany na górze README, nad funkcjami, i wymuszony strukturalną
nieobecnością, a nie polityką:

\begin{quote}
Instrument mierzy \textbf{książkę}, nigdy czytelnika.
\end{quote}

Użyteczne pytanie, na które odpowiada, brzmi: \emph{gdzie ta książka marnuje czas czytelnika},
a nie \emph{który czytelnik jest najgorszy}. Konkretnie, i w formie, którą przeszukanie kodu
może sfalsyfikować: nie ma widoku pojedynczego czytelnika i dodanie go nie jest planowane; nie
ma endpointu, którego trasa albo parametr zapytania nazywa osobę; i nie ma tabeli, z której
dałoby się zbudować ocenę czytelnika. Wyniki są zapisywane przy ramce w danej wersji paczki,
przy podejściu i przy przebiegu sprawdzenia --- przy artefakcie, nigdy przy osobie --- a wiersz
wyniku nie niesie ani identyfikatora, ani znacznika czasu, bo wiersze, których czasy biegną po
kolei przez jeden program, odtwarzają sesję, nikogo nie nazywając.

Pozycja czytelnika w lekturze jest jedyną rzeczą, jaką ten majątek przechowuje o osobie, i
mówi, gdzie ona jest, a nigdy jak jej poszło. Każdą z dwóch tabel trzymają trzy mechaniczne
reguły, a nie obietnica: zamknięta lista kolumn, przez co miara czytelnika psuje build; każdy
klucz i indeks zaczynający się od czytelnika w jednej tabeli, a od wersji paczki w drugiej;
oraz zapytanie, które nie przypina dokładnie jednego z nich, odrzucane w czasie działania,
zanim ORM je skompiluje. Każdą widziano, jak czegoś odmawia, zanim w nią uwierzono.

Czytelnik słyszy także, co go ta nieobecność kosztuje --- a tego strukturalna nieobecność nie
potrafi dostarczyć sama: ponieważ wynik nie niesie czytelnika, nic nie potrafi znaleźć
wierszy, które były jego, więc usunięcie konta nie cofnie wkładu już wliczonego do wskaźnika.
Ekran usuwania mówi dokładnie to, obok tego, co usunięcie faktycznie usuwa.

Dwie dalsze odmowy są zapisane jako decyzje, a nie zostawione jako preferencje. Nigdzie w
pętli czytelnika nie jest wołany żaden model językowy: nauczaniem jest porównanie z kolejną
ramką, a osąd modelu wniósłby własny wskaźnik błędu w materiał dowodowy używany do poprawiania
książki. I nie ma grywalizacji, bo seria jest powodem, by ruszyć liczbę, która nie jest
uczeniem się.

\section{Stan: nic nie jest wdrożone}

Repozytorium stoi na \emph{buduje się, testy zielone, obrazy się budują}. Nie ma działającej
instancji, pod żadnym adresem, dla nikogo. Cztery pliki \code{fly.toml} opisują topologię,
której nigdy nie zastosowano, a żaden z obrazów kontenerów nigdy nie został opublikowany.

Istnieje za to powierzchnia lektury nad całą książką: każdy program w obu edycjach, jedna
ramka na ekran, z własną prozą i matematyką książki; wiersz miejsca, który jest jedyną obudową
i którego numer ramki jest skokiem; oraz arkusz na każdej ramce, która o coś prosi. Za tym:
API, aplikacja webowa, której strona startowa jest indeksem programów, pakiet akceptacyjny
Playwrighta i bramki CI, które złapałyby regresję w którymkolwiek z tych elementów.

Model domeny jest mały i przyszedł późno celowo. Encje wymyślone przed ticketem, który ich
potrzebuje, są kodem, który pierwszy prawdziwy ticket kasuje, więc kontekst trwałości nie
trzymał żadnej, dopóki synchronizacja nie potrzebowała jednej. Dziś trzyma pozycję czytelnika w
lekturze i zliczenie wyników przy ramce --- i nic więcej.

\section{Kształt repozytorium}

\begin{table}[ht]
\centering
\small
\begin{tabular}{@{}ll@{}}
\toprule
\textbf{Ścieżka} & \textbf{Czym jest} \\
\midrule
\code{src/AbOvo.AppHost}        & Korzeń kompozycji. Tylko deweloperski; nie jest topologią produkcyjną. \\
\code{src/AbOvo.ServiceDefaults} & Wspólne jądro: instalacja przekrojowa, pod pułapem rozmiaru. \\
\code{src/AbOvo.Contracts}      & DTO przekraczające granicę serwisu. \\
\code{src/AbOvo.Api}            & Serwis HTTP. Posiada jedną logiczną bazę. Weryfikuje tokeny; nie wydaje żadnych. \\
\code{tests/AbOvo.Api.Tests}    & Jednostkowe i integracyjne w pamięci, plus reguły architektury. Bez kontenera. \\
\code{tests/e2e}                & Pakiet akceptacyjny Playwrighta; własny pakiet i własny lockfile. \\
\code{web/}                     & Przestrzeń pnpm. Jedna aplikacja Next.js i jej backend-for-frontend. \\
\code{flyio/}                   & Wdrożona topologia, opisana i nigdy niezastosowana. \\
\code{docs/}                    & Architektura, decyzje, UX, diagramy i ten papier. \\
\bottomrule
\end{tabular}
\end{table}

Powierzchnia HTTP jest zorganizowana według autoryzacji, a nie według funkcji, i korzeń
kompozycji czyta się jak jej manifest. Grupa publiczna niesie gotowość, żywotność i endpoint
opisujący wdrożenie. Grupa uwierzytelniona niesie własną pozycję czytelnika w lekturze, w tym
jej usunięcie. Grupa administracyjna niesie wskaźniki, które czyta widok autora. Czwarta grupa
nie jest czwartym członkiem tej triady: niesie jeden zapis anonimowy --- wynik --- bo zgoda nie
jest kontem, a instrument wymagający tokenu mierzyłby książkę doświadczaną przez posiadaczy
kont i nazywał to książką. Bierze jawny limit tempa, który bierze grupa uwierzytelniona,
zamiast spadać do limitu globalnego, bo sonda, która dostanie odmowę z powodu limitu, wypada z
rotacji razem z maszyną.

\begin{figure}[ht]
\centering
\includegraphics[width=0.92\linewidth]{../diagrams/rendered/a3-one-origin.pl.pdf}
\caption{Każde żądanie, które wolno wykonać przeglądarce. Dwie strzałki narysowane dlatego, że
nie wolno im istnieć, są tu regułą: przeglądarka rozmawia z własnym originem aplikacji webowej
i z niczym więcej, i dlatego po stronie frontendu nie konfiguruje się CORS.}
\end{figure}

\section{Architektura i to, gdzie odchodzi od konstytucji}

Repozytorium jest budowane wobec architektury referencyjnej trzymanej w osobnym repozytorium
standardów, a \code{docs/architecture/00-ARCHITECTURE.md} przechodzi jej piętnaście zasad po
tym drzewie --- odwołując się do konstytucji, a nie kopiując jej, bo skopiowana zasada jest
drugą kopią, która się dezaktualizuje.

Trzy własności warto nazwać w przeglądzie. Opcjonalne integracje \textbf{degradują się
widocznie}: wdrożenie bez żadnych poświadczeń startuje, serwuje i nazywa to, czego nie ma, w
trzech renderowaniach produkowanych z jednego źródła. Wspólne jądro jest trzymane przy
instalacji przez dwa mechanizmy, a nie przez prozę --- bramkę rozmiaru w CI i test
architektury sprawdzający, że nie referuje ono żadnego typu serwisu i nie deklaruje modelu
trwałości. I żaden adres nie jest wkompilowany w żaden artefakt: frontend czyta swoją
konfigurację w czasie żądania, więc jeden obraz obsługuje każde środowisko.

Ten dokument niesie także \textbf{rejestr odstępstw}, w którym każdy wiersz ma datę, powód i
warunek wyjścia. Co jest w nim zapisane dzisiaj: automatyka wersji zależności zadeklarowana i
wyłączona; bramka rozmiaru licząca wiersze kodu, a nie surowe, żeby zasada wymagająca
uzasadnienia w komentarzach nie wydawała budżetu zasady ograniczającej rozmiar; Python w
przeglądarce, czego standardy nie popierają dowodem; nazwa repozytorium będąca literówką i
nazwa systemu, która nią nie jest; jedno API Aspire oznaczone jako wyłącznie ewaluacyjne,
użyte w pliku wyłącznie deweloperskim; skrypt lustrzany skanu sekretów, który nie jest jeden
do jednego ze swoim jobem CI; oraz paczka treści kompilowana tutaj z książki na przypiętej
rewizji, bo żadne wydanie książki jeszcze takiej nie niesie. Ten ostatni warunek wyjścia jest
najprostszy w całym rejestrze: pierwsze wydanie, które poniesie paczkę.

Odstępstwo bez warunku wyjścia jest dryfem, a nie decyzją, i dlatego ta kolumna istnieje.

\section{Plan}

Pięć faz, uszeregowanych w \code{docs/ux/UI-UX.md} tak, by sesja dostarczająca podnosiła
decyzję, a nie wyprowadzała jej na nowo. Wcześniejsze fazy wylądowały; lista jest trzymana w
pierwotnej kolejności, bo uzasadnienie tej kolejności jest w niej częścią użyteczną.

\begin{enumerate}[leftmargin=2em,itemsep=2pt]
  \item \textbf{Panel laboratorium.} Ćwiczenia, w przeglądarce. Największy kawałek, który nie
        potrzebuje ani backendu, ani paczki treści, i dlatego jest pierwszy.
  \item \textbf{Schemat treści i widok ramki.} Jego druga połowa była \emph{zablokowana} na
        opublikowaniu paczki treści przez książkę --- jedynej zewnętrznej zależności planu ---
        i została odblokowana przez skompilowanie paczki tutaj, na przypiętej rewizji, własnym
        kompilatorem książki i jej własną kontrolą krzyżową. To jest odstępstwo z warunkiem
        wyjścia, a nie rozwiązanie tej zależności. Fikstura napisana na pierwszą połowę
        zostaje jako kontrola warstwy jednostkowej: test, którego wejściem jest badana rzecz,
        dowodzi mniej.
  \item \textbf{Pozycja w lekturze i konta.} Najpierw lokalnie; konto to synchronizuje.
  \item \textbf{Instrument.} Zgoda dobrowolna i wersjonowana; wyniki przy artefaktach; każdy
        wskaźnik niosący swój przedział.
  \item \textbf{Edycja otwartoźródłowa.} Jednorazowa, nieodwracalna bramka, uporządkowana
        według nieodwracalności, w której pełny skan historii pod kątem sekretów blokuje każdą
        inną pozycję.
\end{enumerate}

\section{Gdzie są dokumenty}

\begin{description}[leftmargin=0pt,itemsep=1pt]
  \item[\code{docs/START-HERE.pl.md}] Drzwi frontowe, po polsku i po angielsku: którego z
        czterech rodzajów dokumentu potrzebujesz i gdzie on jest.
  \item[\code{README.md}] Czym system jest, jak uruchomić go jednym poleceniem, i antycele,
        wypowiedziane przed funkcjami.
  \item[\code{docs/architecture/00-ARCHITECTURE.md}] Przejście zasada po zasadzie, rejestr
        odstępstw i znane luki.
  \item[\code{docs/adr/}] Decyzje, po jednym pliku na każdą, wraz z tym, co kosztowały.
  \item[\code{docs/DIAGRAMS.pl.md}] Cały system jako diagramy, w czterech częściach, po polsku
        i po angielsku. Każdy istnieje w treści dokumentu --- jedyna forma, którą renderuje
        GitHub --- oraz jako samodzielny \code{.mmd}, z którego renderowane są rysunki tego
        papieru.
  \item[\code{docs/SCREENSHOTS.pl.md}] Produkt ekran po ekranie, po polsku i po angielsku,
        uchwycony z buildu produkcyjnego przez pakiet akceptacyjny.
  \item[\code{docs/tutorials/}, \code{docs/how-to/}] Dwujęzyczne w całości: lekcje na pierwsze
        uruchomienie i przepisy na zadanie już zrozumiane.
  \item[\code{docs/ux/UI-UX.md}] Ekrany tak, jak zbudowane, i uszeregowany backlog.
  \item[\code{AGENTS.md}, \code{CONTRIBUTING.md}, \code{SECURITY.md}] Praca w repozytorium i
        to, co robić, gdy sekret wyląduje w historii.
\end{description}

\vspace{1em}
\noindent\rule{\linewidth}{0.4pt}

\small
\noindent
Ten dokument nie zapisuje żadnego własnego pomiaru i nie stawia twierdzenia, którego nie
stawia już markdown repozytorium. Jest renderowany na żądanie ręcznym workflow, a powstający
PDF nie jest commitowany.

\end{document}
